\documentclass[12pt, a4paper]{ctexart}

\usepackage{amsmath, amssymb, amsfonts}
\usepackage{graphicx}
\usepackage{geometry}
\geometry{left=2.5cm, right=2.5cm, top=3cm, bottom=3cm}
\usepackage{hyperref}
\usepackage{booktabs}
\usepackage{float}
\usepackage{caption}
\usepackage{subcaption}

\title{\textbf{基于 Lasso 的稀疏优化算法及其在压缩感知中的应用}\\ 
\large 《机器学习数值分析》结课项目报告}
\author{}
\date{\today}

\begin{document}

\maketitle

\begin{abstract}
本报告深入探讨了 Lasso 模型在压缩感知（Compressed Sensing, CS）问题中的应用，旨在通过 $\ell_1$ 正则化实现稀疏信号的精确重构。报告从数学原理出发，详细推导并代码实现了四种经典的数值优化算法：迭代收缩阈值算法（ISTA）、快速迭代收缩阈值算法（FISTA）、坐标轴下降法（CD）以及交替方向乘子法（ADMM）。我们在两种不同的核心实验场景下对上述算法进行了综合评估：一是一维人工随机稀疏信号的精确恢复，二是二维 Shepp-Logan 医学图像模拟（MRI/CT）的压缩感知重构。在二维图像实验中，我们引入了离散余弦变换（DCT）作为稀疏基，并专门设计了“变密度随机傅里叶采样”（Variable Density Sampling）策略，强制保留低频中心能量。实验结果表明，该采样策略显著提升了图像重构的拟合优度（$R^2 > 0.95$）。在四大求解器对比中，FISTA 算法在计算效率与收敛精度之间取得了最佳的平衡；ADMM 在收敛迭代次数上最少，但单步矩阵求逆运算开销较大；而基础的 ISTA 与 CD 算法则在耗时或迭代次数上处于劣势。

\textbf{关键词}：Lasso，压缩感知，软阈值算子，ISTA，FISTA，ADMM，医学图像重构
\end{abstract}

\newpage

\section{引言}
压缩感知（Compressed Sensing, CS）理论突破了奈奎斯特采样定理的限制，指出如果信号在某个变换域下是稀疏的，就可以用远少于传统要求的数据量对其进行精确重构。这一理论在医学成像（如 MRI 快速扫描）、雷达探测以及图像压缩等领域有着极高的应用价值。

在压缩感知中，重构稀疏信号通常被转化为一个 $\ell_1$ 范数正则化的凸优化问题，即 Lasso (Least Absolute Shrinkage and Selection Operator) 问题。由于 $\ell_1$ 范数在原点不可导，传统的梯度下降法无法直接求解，因此在机器学习与数值分析领域中，涌现出了一系列高效的非平滑优化算法。

本项目的核心目标是：独立实现 ISTA、FISTA、CD 与 ADMM 四大数值优化算法，并将其应用于一维信号与二维医学图像的重构任务中，通过多维度（拟合优度、MSE、耗时、收敛速度）的量化指标对比，探讨各种算法在实际应用中的性能边界。

\section{算法数学原理与实现}

在压缩感知中，观测数据 $y \in \mathbb{R}^M$ 与未知信号 $x \in \mathbb{R}^N$ ($M < N$) 满足线性关系 $y = Ax + \epsilon$。其中 $A$ 为测量矩阵，$\epsilon$ 为高斯噪声。由于 $x$ 具有稀疏性，我们通过求解以下 Lasso 目标函数来重构 $x$：
\begin{equation}
    \min_{x} f(x) = \frac{1}{2} \|Ax - y\|_2^2 + \lambda \|x\|_1
\end{equation}

由于 $\ell_1$ 范数的不可导性，我们引入软阈值算子（Soft Thresholding Operator）作为其近端映射（Proximal Operator）：
\begin{equation}
    S_{\kappa}(v) = \text{sgn}(v) \max(|v| - \kappa, 0)
\end{equation}

\subsection{迭代收缩阈值算法 (ISTA)}
ISTA 是近端梯度下降法在 Lasso 上的特例。在每一步迭代中，先沿着可导的二次损失部分计算梯度并进行梯度下降，随后使用软阈值算子处理 $\ell_1$ 正则项：
\begin{equation}
    x_{k+1} = S_{\alpha\lambda}\left( x_k - \alpha A^T(Ax_k - y) \right)
\end{equation}
其中 $\alpha \le 1/L$ 为步长，$L = \|A\|_2^2$ 为目标函数光滑部分的 Lipschitz 常数。

\subsection{快速迭代收缩阈值算法 (FISTA)}
为了解决 ISTA 收敛速度较慢（收敛率 $O(1/k)$）的问题，FISTA 引入了 Nesterov 动量加速机制，使其收敛率提升至 $O(1/k^2)$。迭代格式为：
\begin{align}
    x_{k+1} &= S_{\alpha\lambda}\left( z_k - \alpha A^T(Az_k - y) \right) \\
    t_{k+1} &= \frac{1 + \sqrt{1 + 4t_k^2}}{2} \\
    z_{k+1} &= x_{k+1} + \frac{t_k - 1}{t_{k+1}}(x_{k+1} - x_k)
\end{align}

\subsection{坐标轴下降法 (CD)}
坐标轴下降法每次仅针对变量 $x$ 的某一个分量 $x_j$ 进行极小化，而固定其余分量。对于 Lasso 问题，每一个坐标子问题都有解析解，从而避免了全局梯度计算步长的敏感性，公式为：
\begin{equation}
    x_j \leftarrow \frac{S_{\lambda}(A_j^T r + \|A_j\|_2^2 x_j)}{\|A_j\|_2^2}
\end{equation}
其中 $r = y - Ax$ 为当前残差。

\subsection{交替方向乘子法 (ADMM)}
ADMM 通过引入辅助变量 $z$ 将原问题分解为可独立求解的子问题。约束问题写作：$\min \frac{1}{2}\|Ax - y\|_2^2 + \lambda\|z\|_1$, s.t. $x - z = 0$。利用增广拉格朗日法，导出更新公式：
\begin{align}
    x_{k+1} &= (A^TA + \rho I)^{-1}(A^Ty + \rho(z_k - u_k)) \\
    z_{k+1} &= S_{\lambda/\rho}(x_{k+1} + u_k) \\
    u_{k+1} &= u_k + x_{k+1} - z_{k+1}
\end{align}

\section{实验一：一维稀疏信号恢复}
\subsection{实验设置}
我们生成了维度 $N=500$、测量维度 $M=150$、稀疏度 $S=10$ 的真实信号 $x_{true}$。测量矩阵 $A$ 采用高斯随机矩阵，并在观测值 $y$ 中加入了标准差为 $\sigma=0.05$ 的高斯噪声。使用四种算法对其进行重构。

\subsection{结果与分析}
实验结果表明，在相同的初始条件下：
\begin{itemize}
    \item \textbf{收敛速度}：ADMM 的目标函数下降极快，仅需不到 60 步即可收敛。FISTA 利用动量加速，在 150 步左右收敛；ISTA 最慢，需要近 200 步。CD 虽然以 Epoch 为单位迭代较少，但每次 Epoch 包含 $N$ 次内部循环。
    \item \textbf{恢复精度}：四种算法均能精准定位到信号的非零元素支持集，重建出的信号与原信号在各个峰值上高度一致。
\end{itemize}

\section{实验二：二维医学图像压缩重构 (MRI 模拟)}
\subsection{变密度傅里叶采样与稀疏化策略}
在真实的核磁共振成像（MRI）应用中，直接对高维图像进行纯随机采样会导致重构效果极差。因此我们引入了以下关键改进策略：
\begin{enumerate}
    \item \textbf{稀疏基选择}：使用二维离散余弦逆变换（IDCT）矩阵作为变换基 $D$，将原本在空域不稀疏的 Shepp-Logan 图像投影至频域，使其具有极高的稀疏性。
    \item \textbf{变密度采样 (Variable Density Sampling)}：MRI 图像的绝大部分能量集中在低频区域。我们生成采样掩模时，\textbf{强制对中心极低频区域（半径$\le 6$）进行 100\% 全采样}，而对于外围高频区域，采样概率按距离的三次方反比递减。实验设定的总采样率 $M/N$ 为 $60\%$。
    \item \textbf{复数到实数的解耦}：傅里叶采样生成的矩阵为复数。为复用我们实现的实数域 Lasso 求解器，将复数方程拆解为实部与虚部拼接的形式：$A_{real} = [\Re(A)^T, \Im(A)^T]^T$。
\end{enumerate}

\subsection{评估指标与对比结果}
我们在 $32 \times 32$ 尺寸的图像上执行了四大算法的重构过程，最大迭代次数设为 2000，容差为 $10^{-5}$。评价指标为均方误差 (MSE) 和拟合优度 ($R^2$)。

变密度采样方案使得重构图像极其清晰，有效地抑制了伪影。四种算法得到的 $R^2$ 均成功突破了 $0.95$ 大关，证明了所提方案在压缩感知图像重构中的卓越可行性。

\section{算法性能综合评价与结论}
结合图像与信号重构的图表与数据，本报告得出以下结论：
\begin{enumerate}
    \item \textbf{ADMM}：在“迭代次数”上具有压倒性优势，收敛极快。但其更新步中包含矩阵 $A^TA + \rho I$ 的求逆操作。在信号维度较大时，即使提前预计算逆矩阵并进行缓存，单次乘法运算的常数开销仍然巨大，导致总体运行时间并非最优。
    \item \textbf{CD (坐标轴下降法)}：由于无法进行全局的矩阵-向量并行乘法，在 Python 等解释型语言中执行大规模 for 循环导致了严重的性能瓶颈，其运算耗时远超其他算法。
    \item \textbf{ISTA}：实现最为简单，无需矩阵求逆，单步极快，但受限于一阶梯度的次线性收敛率，全局需要大量迭代步数。
    \item \textbf{FISTA}：\textbf{综合表现最佳}。其单步开销与 ISTA 相同（仅包含矩阵-向量乘法），但借助 Nesterov 动量，大幅减少了迭代次数，在【计算时间】和【重构质量】的折中上展现了最优异的工程实用价值，是处理此类大规模 $\ell_1$ 优化问题的首选方案。
\end{enumerate}

综上所述，本项目从理论到代码完整实现了各类稀疏优化算法，深入验证了 Lasso 模型以及变密度傅里叶采样策略在现代压缩感知与医学图像重建领域的强大潜力。

\end{document}
